% Prostudováno: ucebnice/02_Funkce_a_rovnice/Linearni_funkce_a_rovnice/Linerání rovnice.tex,
% část Lineární rovnice o jedné neznámé, Ekvivalentní úpravy, Zkouška (ř. 57–310).
% Realisticky.cz: externi_zdroje/realisticky/02_Funkce_a_rovnice/02_Linearni_rovnice_a_nerovnice/020201_Rovnice_ekvivalentni_upravy/020201_Lekce.pdf
% a 020203_Linearni_rovnice_I/020203_Lekce.pdf; použita myšlenka stejných operací na obou stranách a vazba na průsečík grafu. Úlohy vlastní.
\documentclass[11pt,a4paper]{article}
% Spolecny styl pro ucitelske pripravy MAT (Work/MAT)
% Verze 1.1, 2026-10-23
\NeedsTeXFormat{LaTeX2e}
\ProvidesPackage{priprava}[2026/10/23 v1.1 Sablona ucitelskych priprav MAT]

% Zakladni typografie a cestina
\RequirePackage[T1]{fontenc}
\RequirePackage[utf8]{inputenc}
\RequirePackage[czech]{babel}
\RequirePackage{lmodern}
\RequirePackage{microtype}
\RequirePackage{amsmath,amssymb}
\RequirePackage{xcolor}
\RequirePackage{enumitem}
\RequirePackage[most]{tcolorbox}
\RequirePackage{tabularx,booktabs,array}
\RequirePackage{fancyhdr}
\RequirePackage{needspace}
\RequirePackage{geometry}
\RequirePackage[hidelinks]{hyperref}

\geometry{a4paper,top=17mm,bottom=18mm,left=23mm,right=23mm}
\definecolor{notegray}{gray}{0.39}
\setlength{\parindent}{0pt}
\setlength{\parskip}{2.5pt plus 1pt minus 0.5pt}
\setlist[itemize]{topsep=2pt,itemsep=1pt,parsep=0pt,leftmargin=1.45em}
\setlist[enumerate]{topsep=2pt,itemsep=2pt,parsep=0pt,leftmargin=1.65em}
\setcounter{secnumdepth}{0}
\renewcommand{\arraystretch}{1.12}

% Zahlavi: vlevo rocnik | globalni tema, vpravo cislo: nazev hodiny
% Zapati zustava zcela prazdne.
\pagestyle{fancy}
\fancyhf{}
\renewcommand{\headrulewidth}{0.35pt}
\renewcommand{\footrulewidth}{0pt}
\setlength{\headheight}{14pt}
\newcommand{\nastavPripravu}[4]{%
  \fancyhead[L]{\small #1\enspace\textbar\enspace #2}%
  \fancyhead[R]{\small #3: #4}%
}

% Minimalisticke cernobile boxy. Kazda definice zvlast.
\tcbset{commonlesson/.style={enhanced,breakable=false,colback=white,colframe=black,
  colbacktitle=white,coltitle=black,boxrule=0.45pt,
  arc=0.8mm,outer arc=0.8mm,left=3mm,right=3mm,top=1.6mm,bottom=1.6mm,
  toptitle=1mm,bottomtitle=0.7mm,fonttitle=\bfseries,
  attach boxed title to top left={xshift=2.5mm,yshift=-1.15mm},
  boxed title style={colback=white,colframe=white,boxrule=0pt,left=1mm,right=1mm}}}
\newtcolorbox{zadaniBox}[1][]{commonlesson,title={Zadání},#1}
\newtcolorbox{prikladBox}[1][]{commonlesson,title={Příklad},#1}
\newtcolorbox{definiceBox}[1][]{commonlesson,title={Definice},#1}

% Spolecny box pro zadani a jeho reseni, bez mezery mezi castmi.
% Pouziti:
%   \begin{prikladReseniBox}
%     Zadani ulohy.
%     \tcblower
%     \begin{align*} ... \end{align*}
%   \end{prikladReseniBox}
% Spodni cast ma jemny sedy podklad a delici linku; lze v ni sazet
% vice radku matematiky. Obe casti maji jen jeden spolecny ramecek.
\tcbset{solutionlesson/.style={
  commonlesson,
  enhanced jigsaw,
  breakable,
  colbacklower=black!3,
  segmentation style={solid,draw=black!25,line width=0.35pt},
  fontlower=\small\color{notegray},
  before lower={\noindent\textbf{Řešení:}\enspace},
  leftlower=3mm,rightlower=3mm
}}
\newtcolorbox{prikladReseniBox}[1][]{solutionlesson,title={Příklad},#1}
\newtcolorbox{zadaniReseniBox}[1][]{solutionlesson,title={Zadání},#1}

% Pomucky pro autora pripravy: sedive poznamky; \reseni pro starsi pripravy
\newcommand{\poznamka}[1]{\par\begingroup\small\color{notegray}\itshape
  \textnormal{\textbf{Poznámka pro učitele.}}\enspace #1\par\endgroup}
\newcommand{\reseni}[1]{\par\begingroup\small\color{notegray}
  \textbf{Očekávané řešení:}\enspace #1\par\endgroup}
\newcommand{\blok}[1]{\needspace{5\baselineskip}\subsection*{#1}}
\newcommand{\vyrok}[1]{\uv{#1}}

\usepackage{tikz}
\nastavPripravu{1. ročník}{Rovnice a nerovnice}{73}{Lineární rovnice a ekvivalentní úpravy}
\begin{document}
\blok{Rovnice jako hledání argumentu}
V hodině 72 jsme hledali, kdy mají dvě funkce stejnou hodnotu. Nyní chceme určit všechny argumenty, které rovnost splňují, a přitom umět zdůvodnit každý krok. Pro výrazy $L(x)$ a $P(x)$ může jeden dosazený argument rovnost splnit, jiný ne.
\begin{definiceBox}
\textbf{Kořen rovnice} $L(x)=P(x)$ je číslo $x_0$ z přípustného oboru, pro které $L(x_0)=P(x_0)$. Množinu všech kořenů označíme $K$.
\end{definiceBox}
\begin{definiceBox}
\textbf{Lineární rovnice o jedné neznámé} je rovnice převoditelná na $ax+b=0$, kde $a,b\in\mathbb R$ a $a\ne0$. Při $a=0$ dostáváme mezní případ $b=0$, který vyhodnotíme zvlášť jako pravdivý či nepravdivý výrok.
\end{definiceBox}
\begin{definiceBox}
\textbf{Ekvivalentní úprava rovnice} je úprava, která zachovává množinu všech řešení. Na obou stranách lze přičíst či odečíst stejný výraz; násobit nebo dělit lze stejným číslem různým od nuly. U výrazu s neznámou je nutno před dělením ověřit nenulovost pro všechny uvažované hodnoty nebo rozdělit případy.
\end{definiceBox}
\begin{definiceBox}
\textbf{Zkouška} nalezeného kořene spočívá v jeho dosazení zvlášť do levé a pravé strany původní rovnice. Shodné hodnoty potvrzují, že kandidát rovnici splňuje; úplnost množiny řešení musí plynout z postupu.
\end{definiceBox}
\begin{prikladReseniBox}
Řešte $3x-5=x+7$ a proveďte zkoušku. Každou změnu vysvětlete jako operaci s oběma stranami.
\tcblower
Odečteme $x$: $2x-5=7$. Přičteme $5$: $2x=12$. Vydělíme nenulovým číslem $2$: $x=6$, tedy $K=\{6\}$. Zkouška v původní rovnici: $L(6)=18-5=13$, $P(6)=6+7=13$. Všechny úpravy byly ekvivalentní, takže jiné řešení není.
\end{prikladReseniBox}
\begin{center}
\begin{tikzpicture}[x=.65cm,y=.5cm,>=stealth]
\draw[step=1,gray!25,very thin] (3,-4) grid (8,4);
\draw[->] (2.7,0)--(8.4,0) node[right] {$x$};
\foreach \x in {4,5,6,7,8}{\draw (\x,.1)--(\x,-.1) node[below] {\small $\x$};}
\draw[very thick] (4,-4)--(8,4);
\fill (6,0) circle (2.6pt) node[below right=6pt] {\small $(6;0)$};
\end{tikzpicture}
\end{center}
\poznamka{Graf rozdílu $y=(3x-5)-(x+7)=2x-12$ protíná osu $x$ právě při řešení původní rovnice. Jde o přesný graf v zobrazeném okně; hodnoty obou stran rovnice při $x=6$ jsou shodně 13.}
\poznamka{Jádro: motivace a definice 12 min, společné příklady a zkouška 15 min, samostatná práce s kontrolou 15 min, organizace 3 min. Rezerva nad rámec, domácí práce asi 15 min. Další hodina přidá závorky a zlomky.}
\newpage
\blok{Rovnice s více či žádným řešením}
\begin{prikladReseniBox}[breakable=false]
Rozhodněte o množině řešení rovnic $2(x+3)=2x+6$ a $2(x+3)=2x+5$.
\tcblower
První po roznásobení dává $2x+6=2x+6$, tedy $0=0$ pro každé reálné $x$ a $K=\mathbb R$. Druhá dává $2x+6=2x+5$, tedy $6=5$, což je nepravda pro každé $x$; $K=\varnothing$. Není dovoleno z obou rovnic po zrušení $x$ usoudit stejný výsledek.
\end{prikladReseniBox}
\begin{prikladReseniBox}[breakable=false]
Proč není dělení výrazem $x-2$ bez další úvahy ekvivalentní u rovnice $(x-2)(x+1)=0$?
\tcblower
Původní rovnice má kořeny $x=2$ a $x=-1$. Dělení $x-2$ by pro $x=2$ znamenalo dělení nulou a tento kořen by ztratilo. Bezpečný postup je nulový součin: $x-2=0$ nebo $x+1=0$; kontrola obou kořenů v původním součinu dává nulu.
\end{prikladReseniBox}
\blok{Samostatné procvičení}
\begin{zadaniReseniBox}
\begin{enumerate}[label=\alph*)]
\item Řešte $5x+4=2x-8$ a ověřte kořen.
\item Řešte $4(x-1)=2(2x+1)$; určete množinu řešení.
\item Žák z $-3x=12$ odvodil $x=4$. Najděte chybu a ověřte opravu.
\end{enumerate}
\tcblower
\begin{enumerate}[label=\alph*)]
\item $3x=-12$, $x=-4$. Zkouška: $5(-4)+4=-16$ a $2(-4)-8=-16$; $K=\{-4\}$.
\item $4x-4=4x+2$ vede k $-4=2$, tedy $K=\varnothing$.
\item Dělení číslem $-3$ dává $x=-4$. Kontrola: $-3(-4)=12$; $x=4$ by dalo $-12\ne12$.
\end{enumerate}
\end{zadaniReseniBox}
\blok{Volitelné rozšíření a rezerva}
\begin{zadaniReseniBox}
U rovnice $a x+b=0$ pro daná reálná čísla $a,b$ rozlište případy $a\ne0$, $a=0,b\ne0$ a $a=b=0$.
\tcblower
Pro $a\ne0$ je jediný kořen $x=-b/a$. Pro $a=0,b\ne0$ je rovnice $b=0$ nepravdivá, tedy $K=\varnothing$. Pro $a=b=0$ je $0=0$ vždy pravdivé, tedy $K=\mathbb R$. Nejde zatím o obecný rozbor parametru v zadání; ten přijde v hodině 76.
\end{zadaniReseniBox}
\blok{Domácí práce}
\begin{zadaniReseniBox}
\begin{enumerate}[label=\alph*)]
\item Vyřešte $7-2x=x+1$ a proveďte zkoušku.
\item Rozhodněte, zda $3(x+2)-x=2x+6$ má jeden, žádný či nekonečně mnoho kořenů.
\item Vysvětlete na příkladu, proč dělení rovnice nulou není přípustné.
\end{enumerate}
\tcblower
\begin{enumerate}[label=\alph*)]
\item $6=3x$, $x=2$; $L=7-4=3$, $P=2+1=3$, tedy $K=\{2\}$.
\item Po úpravě je $2x+6=2x+6$, tedy pravda pro všechna $x\in\mathbb R$ a $K=\mathbb R$.
\item Např. $2x=4$ má $x=2$. Vynásobení obou stran nulou dává $0=0$ pro každé $x$ a ztratí informaci o původním kořeni; dělení nulou vůbec není definováno.
\end{enumerate}
\end{zadaniReseniBox}
\end{document}
